\input{./._relpath-to-latexroot.ltx}
\documentclass[\latexroot/\projectname]{subfiles}
\input{\latexroot/@resources/subfile-setup.ltx}

% Model.tex-specific: Redefine verbatimwrite to be a no-op in standalone mode
% (whenstandalone removed - redundant in preamble)
\makeatletter
\renewenvironment{verbatimwrite}[1]{\comment}{\endcomment}
\makeatother

\begin{document}

% When compiled as a subfile of an integrated document, path reset
% is needed so that paths are relative to the root file location
\whenintegrated{%
  \input{\latexroot/._relpath-to-latexroot.ltx}%  % Resets \latexroot and reloads all paths
}

\section{Model}
\whenintegrated{\label{sec:model}} % integrated doc: SST for labels

This section presents our model. It features multiple regions that engage in intermediate-input trade, heterogeneous households, and incomplete markets.

\subsection{Description of the economy}
The model is dynamic and solved in general equilibrium. In the model description below, we abstract from the time subscript $t$. The economy has $N=2$ regions: one (small) region and the

\subfile{../Figures/local-multiplier-trade}

rest of the economy. Nakamura and Steinsson~\cite{Nakamura2011-dw} also assume this geographical representation. To be consistent with our empirical evidence (which is based on county-level responses), we calibrate the size of the small region according to the population size of the average county in the U.S.

Each region $i$ has its own real wage $w_{i}$ and inflation rate $\pi_{i}$. Each region produces a final good $Y_{i}$ using intermediate inputs produced by monopolistically competitive firms. There is a continuum of intermediate good firms in each region, indexed by $(i, j)$. Intermediate good firm $j$ located in region $i$ produces $y_{i, j}$ at price $p_{i, j}$. Regional trade takes place through trade in intermediate inputs. Regional trade flows are calibrated based on data from the CFS that combine shipments of final and intermediate goods. Through the lens of the model, we interpret all trade flows (including tradeable services) as trade in the intermediate sector and assume that the final goods are only locally consumed.\footnote{Our model allows trade in goods but not labour movement between regions. Auerbach et al.~\cite{Auerbach2019-pi} show that local federal spending does not pull in labour from nearby locations. House et al.~\cite{House2020-ho} also explore fiscal policy in a multi-country dynamic stochastic general equilibrium (DSGE) setup with explicit trade linkages. In their model, each country has a representative household.}

The region $i$ population is denoted $\mu_{i}$ with $\sum_{i} \mu_{i}=1$. Populations do not vary across time. The regions are symmetric, i.e. per-capita variables are identical across regions at the steady state. In each region, there is a continuum of households making consumption, working, and saving decisions. Finally, there is a government buying final goods from each region. To finance expenditures, it taxes households' labour income. Taxation occurs only at the federal level (fiscal union). The government also supplies the nominal bond used by households as a savings instrument. Households in both regions face the same nominal interest rate $R$ (currency union). Therefore, there are two sources of market incompleteness in the model: first, across regions, and second, within regions.

Household-level variables are denoted with a small letter. Per-capita regional variables are denoted with capital letters. Aggregate variables are per-capita variables times the population rate. For example, consumption of a household in region $i$ is $c_{i}$, per-capita consumption in region $i$ is $C_{i}$, and aggregate consumption in region $i$ is $\mu_{i} C_{i}$.

\subsection{Households}
Each region is populated by a measure one continuum of households. Households derive utility from consumption (denoted $c$ ) and leisure. The good is fully consumed at the time of the purchase and does not provide any service beyond the current period. Hence, consumption corresponds to non-durable spending. A household is endowed with one unit of productive time, which it splits between work $h$ and leisure. Households' decisions depend on preferences represented by a time separable utility function of the form


\begin{equation*}
U=\mathbf{E}_{0}\left[\sum_{t=0}^{\infty} \beta^{t-1}\left\{\log \left(c_{i, t}\right)+\psi \frac{\left(1-h_{i, t}\right)^{1-\theta}}{1-\theta}\right\}\right], \tag{4.1}
\end{equation*}


where $\psi$ affects the utility from leisure and $\theta$ affects the Frisch elasticity of labour supply.\footnote{Given the utility specification, the Frisch elasticity of labour supply depends on both $\theta$ and the hours worked $h$. In Section 6, we analyse a utility function specified in terms of disultility of hours worked where $\theta$ is directly interpreted as the labour supply elasticity.} $\beta$ is the household's idiosyncratic discount factor which does not vary across time (~\cite{Krueger2016-kh}; ~\cite{Hagedorn2019-op}). Half of the households in each region are impatient ( $\beta= \bar{\beta}-\epsilon_{\beta}$ ) and half are patient ( $\beta=\bar{\beta}+\epsilon_{\beta}$ ). Discount factor heterogeneity allows the model to be consistent with evidence of non-durable consumption responses to tax rebates.

Households consume only the final good produced in their region. They supply labour in the intermediate good sector of their region and receive real wage payments $w_{i}$. Their effective labour supply is $x h$ where $x$ is an idiosyncratic shock that follows an AR(1) process in logs:


\begin{equation*}
\log x_{t+1}=\rho \log x_{t}+\eta_{t+1} \quad \text { with } \eta_{t+1} \sim \operatorname{iid} N\left(0, \sigma_{\eta}^{2}\right) \tag{4.2}
\end{equation*}


Therefore, households are heterogeneous for two reasons: (i) they have different discount factors and (ii) they receive in every period different shocks to their productivity. The differences in patience and labour earnings result in differences in consumption and asset holdings. The transition matrix that approximates the autoregressive process in equation (4.2) is given by $\Gamma_{x x^{\prime}}$.

In each period $t$, a household residing in region $i$ saves $a_{i, t+1}$ in a regional mutual fund. The mutual fund is discussed in detail below. One period later, the household receives from the mutual fund a real return equal to $r_{i, t+1}$. Following Kaplan et al.~\cite{Kaplan2018-sr}, we assume that households receive a $1-\omega$ fraction of real dividends, $D_{i, t}$, directly from the local intermediate firms. This additional income is interpreted as the profit-sharing component of worker compensation such as bonuses, commissions, and stock options. The residual share of dividends $(1-\omega) D_{i, t}$ is distributed proportionately to workers' labour productivity $\delta(x)=x / \bar{x}$ where $\bar{x}$ is the average labour productivity. Finally, households pay income taxes on the sum of wage payments and bonuses, based on the tax schedule $\mathcal{T}($.$) . We denote the regional distribution of households$ across productivity, asset holdings, and discount rates as $\phi_{i}$.

We write the decision problem of a household that resides in region $i$. For simplicity, we index only regional and not idiosyncratic variables by $i$.


\begin{align*}
V_{t}\left(x_{t}, a_{t}, \beta ; \phi_{i, t}\right)= & \max _{c_{t}, a_{t+1}, h_{t}}\left\{\log \left(c_{t}\right)+\psi \frac{\left(1-h_{t}\right)^{1-\theta}}{1-\theta}\right. \\
& \left.+\beta \sum_{x_{t+1}} \Gamma_{x_{t}, x_{t+1}} V_{t+1}\left(x_{t+1}, a_{t+1}, \beta ; \phi_{i, t+1}\right)\right\}  \tag{4.3}\\
\text { s.t. } \quad & c_{t}+a_{t+1}=w_{i, t} x_{t} h_{t}-\mathcal{T}(.)+\left(1+r_{i, t}\right) a_{t}+(1-\omega) \delta\left(x_{t}\right) D_{i, t},  \tag{4.4}\\
& a_{t+1} \geq 0 . \tag{4.5}
\end{align*}


\subsection{Mutual fund}
There is a mutual fund in each region. The fund collects savings from local households, $A_{i, t+1}= \int_{\phi_{i, t}} a_{i, t+1}$. The fund next purchases government bonds, $B_{i, t+1}^{m}$, and shares, $s_{i, t+1}$, that represent claims on the profits of intermediate firms in region $i$. The price of each share (relative to the price of the final good) is $q_{i, t}$. The mutual fund chooses $B_{i, t+1}^{m}$ and $s_{i, t+1}$ to maximize its real long-run profits:


\begin{align*}
V_{i, t}^{m}= & \left(1+R_{t-1}\right) B_{i, t}^{m}+\left(q_{i, t}+\left(1-\tau_{d}\right) \omega D_{i, t}\right) s_{i, t}-\left(1+\pi_{i, t+1}\right) B_{i, t+1}^{m}-q_{i, t} s_{i, t+1} \\
& +\frac{1}{1+r_{i, t+1}} V_{i, t+1}^{m} . \tag{4.6}
\end{align*}


The government bond costs $\$ 1$ and pays $(1+R)$ dollars where $R$ is the nominal interest rate. $\pi_{i, t+1}$ is the region-specific inflation rate defined as


\begin{equation*}
\pi_{i, t+1}=\frac{P_{i, t+1}}{P_{i, t}}-1 \tag{4.7}
\end{equation*}


where $P_{i}$ is the price of the final good in region $i$ (defined below). Each period, the mutual fund receives after-tax dividends $\left(1-\tau_{d}\right) \omega D_{i, t}$, where $\omega$ is the fraction of dividends not distributed directly as employment compensation. The mutual fund pays back households using the revenues from holding government bonds and shares (zero profit condition):


\begin{equation*}
\left(1+r_{i, t}\right) A_{i, t}=\left(1+R_{t-1}\right) B_{i, t}^{m}+\left(q_{i, t}+\left(1-\tau_{d}\right) \omega D_{i, t}\right) s_{i, t} \tag{4.8}
\end{equation*}


Solving the mutual fund's problem gives rise to a no-arbitrage condition:


\begin{equation*}
1+r_{i, t+1}=\frac{1+R_{t}}{1+\pi_{i, t+1}}=\frac{q_{i, t+1}+\left(1-\tau_{d}\right) \omega D_{i, t}}{q_{i, t}} \tag{4.9}
\end{equation*}


Combining equations (4.8) and (4.9) gives the total amount of savings in region $i$ :


\begin{equation*}
A_{i, t+1}=\left(1+\pi_{i, t+1}\right) B_{i, t+1}^{m}+q_{i, t} s_{i, t+1} \tag{4.10}
\end{equation*}


\subsection{Firms}
We describe here the problems of the final good firms and the intermediate good firms. Unless stated otherwise, we omit the time subscript $t$.

Final good firms. There is one final good firm in every region $i$ that produces $Y_{i}$. Each final good is sold at $P_{i}$, which is the price aggregator in each region $i . Q_{i^{\prime}, i}$ is the relative price (real exchange rate) between final goods $i^{\prime}$ and $i$ : $Q_{i^{\prime}, i}=P_{i}^{\prime} / P_{i}$. Each final good uses a variety of intermediate inputs. Inputs are purchased not only locally but from other regions as well. Let the demand from region $i$ of input $j$ that is produced in region $i^{\prime}$ be noted $v_{i, i^{\prime}, j}$. It is purchased at price $p_{i^{\prime}, j}$. The production technology is


\begin{equation*}
Y_{i}=\left[\sum_{i^{\prime}=1}^{N} \gamma_{i i^{\prime}}^{1 / \epsilon} \int_{j} v_{i, i^{\prime}, j}^{(\epsilon-1) / \epsilon} \mathrm{d} j\right]^{\epsilon /(\epsilon-1)} \tag{4.11}
\end{equation*}


Parameter $\gamma_{i i^{\prime}}$ denotes the preference of firm $i$ for inputs from region $i^{\prime}$. We assume that $\sum_{i^{\prime}} \gamma_{i i^{\prime}}=1$. Home bias for region 1 is given by $\gamma_{11}=\alpha$ so that $\gamma_{12}=1-\alpha$. If region 1 imports $1-\alpha$, then region 2 imports $\gamma_{21}=\mu_{1} / \mu_{2} \times(1-\alpha)$, and home bias for region 2 is $\gamma_{22}=1-\left(\mu_{1} / \mu_{2}\right)(1-\alpha)$. The parameter $\epsilon$ captures the substitutability between intermediate inputs. Demand of final good firm $i$ for input $j$ located at $i^{\prime}$ is


\begin{equation*}
v_{i, i^{\prime}, j}=\gamma_{i i^{\prime}}\left[\frac{p_{i^{\prime}, j}}{P_{i}}\right]^{-\epsilon} Y_{i} \tag{4.12}
\end{equation*}


The final good firm is making zero profits (perfect competition), which allows us to write the price aggregate as


\begin{equation*}
P_{i}=\left[\sum_{i^{\prime}=1}^{N} \gamma_{i i^{\prime}} \int_{j} p_{i^{\prime}, j}^{1-\epsilon} \mathrm{d} j\right]^{1 /(1-\epsilon)} \tag{4.13}
\end{equation*}


Intermediate good firms. Each region $i$ has a continuum of intermediate goods indexed by $j$. The intermediate good $y_{i, j}$ is produced using only labour. We assume that labour cannot move across regions. Firms use a linear technology


\begin{equation*}
y_{i, j}=L_{i, j} \tag{4.14}
\end{equation*}


where $L_{i, j}$ is labour demanded by firm $j$ in region $i$. The intermediate good firm faces demand both from the local and the foreign final good firm. As mentioned, firm $j$ located in region $i^{\prime}$ faces demand by final good firm $i$ equal to $v_{i, i^{\prime}, j}$. The aggregate demand for region $i$ intermediate good firm $j$ will be


\begin{equation*}
y_{i, j}=\sum_{i^{\prime}} \mu_{i^{\prime}} v_{i^{\prime}, i, j} \tag{4.15}
\end{equation*}


Due to monopolistic competition, the intermediate good firm takes its demand into account when setting its price $p_{i, j}$. The intermediate good firm discounts the future based on the rate of return $r$. We denote $\tilde{\beta}_{i}=1 /\left(1+r_{i, t+1}\right)$. Each firm can adjust its price with probability $\lambda$. Since the intermediate good firm is solving a dynamic problem, we re-introduce the time subscript $t$ into our equations. The reset price $p^{*}$ is found by maximizing the value of firm


\begin{equation*}
\max _{p_{i, j, t}^{*}} \sum_{s=0}^{\infty}(1-\lambda)^{s} \frac{1}{\Pi_{h=1}^{s}\left(1+r_{t+h}\right)}\left\{p_{i, j, t+s}^{*} y_{i, j, t+s}-W_{i, t+s} L_{i, j, t+s}\right\} \tag{4.16}
\end{equation*}


where $W_{i, t}$ is the nominal wage. This leads to the optimal pricing equation


\begin{equation*}
\frac{p_{i, j, t}^{*}}{P_{i, t}}=\frac{\epsilon}{\epsilon-1} \frac{\sum_{i^{\prime}=1}^{N} \mu_{i^{\prime}} \gamma_{i^{\prime} i} Q_{i^{\prime}, i, t}^{\epsilon}\left[w_{i, t} Y_{i^{\prime}, t}+(1-\lambda) \tilde{\beta}_{i}\left(1+\pi_{i^{\prime}, t+1}\right)^{1+\epsilon} \mathcal{X}_{i^{\prime}, i, t+1}\right]}{\sum_{i^{\prime}=1}^{N} \mu_{i^{\prime}} \gamma_{i^{\prime} i} Q_{i^{\prime}, i, t}^{\epsilon}\left[Y_{i^{\prime}, t}+(1-\lambda) \tilde{\beta}_{i}\left(1+\pi_{i^{\prime}, t+1}\right)^{\epsilon} \mathcal{Z}_{i^{\prime}, t+1}\right]} \tag{4.17}
\end{equation*}


with


\begin{align*}
\mathcal{X}_{i^{\prime}, i, t} & =w_{i, t} Y_{i^{\prime}, t}+(1-\lambda) \tilde{\beta}_{i}\left(1+\pi_{i^{\prime}, t+1}\right)^{1+\epsilon} \mathcal{X}_{i^{\prime}, i, t+1}  \tag{4.18}\\
\mathcal{Z}_{i^{\prime}, i, t} & =Y_{i^{\prime}, t}+(1-\lambda) \tilde{\beta}_{i}\left(1+\pi_{i^{\prime}, t+1}\right)^{\varepsilon} \mathcal{Z}_{i^{\prime}, i, t+1} \tag{4.19}
\end{align*}


Finally, the real profits for intermediate firm $j$ in region $i$ are


\begin{equation*}
d_{i, j, t}=\frac{p_{i, j, t}}{P_{i, t}} y_{i, j, t}-w_{i, t} y_{i, j, t} \tag{4.20}
\end{equation*}


and the per-capita real dividends in region $i$ are $D_{i}=\int_{j} d_{i, j}$.

\subsection{Monetary authority}
Regions are part of a monetary union. We consider a simple Taylor rule where the monetary authority sets the nominal rate based on the aggregate inflation rate $\hat{\pi}$. In particular,


\begin{equation*}
R_{t}=R_{s s}+\zeta \hat{\pi}_{t} \tag{4.21}
\end{equation*}


The aggregate inflation rate $\hat{\pi}=\sum_{i} \mu_{i} \pi_{i, t}$ is a weighted average of the regional inflation rates. Since after 2008, the monetary policy was constrained by the lower bound, we set $\zeta=0$ in our benchmark calibration. This case captures the effect of government spending in an environment where the monetary authority is unresponsive to inflation pressures.

\subsection{Government}
The government buys final goods from every region. Per-capita government spending in region $i$ is denoted $G_{i}$. It finances spending using labour income taxes and dividend taxes. Labour income taxes are the sum of a lump-sum component $T$ and a proportional tax $\tau$ :

$$
\mathcal{T}_{t}=-T+\tau\left[w_{i t} x_{t} h_{t}+\delta\left(x_{t}\right)(1-\omega) D_{i t}\right]
$$

The government budget constraint is


\begin{equation*}
\sum_{i} \mu_{i}\left(1+\pi_{i, t+1}\right) B_{i, t+1}^{m}-\left(1+R_{t-1}\right) \sum_{i} \mu_{i} B_{i, t}^{m}=\sum_{i} \mu_{i} G_{i, t}-\sum_{i} \mu_{i} \int_{\phi_{i, t}} \mathcal{T}_{t}-\tau_{d} \omega \sum_{i} \mu_{i} D_{i t}, \tag{4.22}
\end{equation*}


where $B_{i, t}^{m}$ is the bond holdings of the regional mutual fund. In equilibrium, the demand for government bonds by the regional mutual funds is equal to the supply of government bonds by the government: $\sum_{i} \mu_{i} B_{i, t}^{m}=B_{t}^{g}$. Along the transition, the supply of bonds is given by the following fiscal rule: $B_{t}^{g}=B_{s s}^{g}-\sum_{i} \mu_{i}\left(q_{i, t}-q_{i, s s}\right)$. According to this fiscal rule, any decline in equity value is met with an equivalent increase in government debt. As is typical in models with sticky prices, government spending increases the nominal wages, reducing the markup and therefore the firms' dividends and the value of equity. The equity shares are implicitly held\\
by households (together with government bonds). We neutralize the effect of movements in equity value by making the assumption that the government intervenes and injects liquidity in the economy in the form of an additional supply of government bonds.\footnote{In Kaplan et al.~\cite{Kaplan2018-sr}, the equity is part of illiquid assets. As a result, they can neutralize the effects of a countercyclical movement in equity value by assuming a profit distribution rule that guarantees an illiquid income flow that is independent of the markup. We cannot rely on a similar assumption as our model features a single asset.}

\subsection{Regional accounts}
We describe the regional income accounts once more abstracting from time subscript $t$. Regional income is equal to the total value added by all intermediate firms in that region: $\mu_{i} \mathcal{Y}_{i}= \int_{j} \sum_{i^{\prime}} \mu_{i^{\prime}} v_{i^{\prime}, i, j} \mathrm{~d} j$. Per-capita income for every region $i$ is equal to $\mathcal{Y}_{i}=w_{i} L_{i}+D_{i}$. Per-capita final $\operatorname{good} Y_{i}$ is equal to per-capita consumption $C_{i}$ plus per-capita government spending $G_{i}$.

\subsection{Characterizing the equilibrium}
We derive expressions that clarify some of our equilibrium conditions. As mentioned, the total demand for intermediate firm $(i, j)$ in period $t$ is

$$
y_{i, j, t}=\sum_{i^{\prime}} \mu_{i^{\prime}} v_{i^{\prime}, i, j, t}=\sum_{i^{\prime}} \mu_{i^{\prime}} \gamma_{i^{\prime} i}\left[\frac{p_{i, j, t}}{P_{i^{\prime}, t}}\right]^{-\epsilon} Y_{i^{\prime}, t}
$$

Aggregating over $j$, we derive the total demand for intermediate inputs of region $i$ in period $t$


\begin{align*}
\int_{j} y_{i, j, t}=\mu_{i} \mathcal{Y}_{i, t} & =\sum_{i^{\prime}} \mu_{i^{\prime}} \gamma_{i^{\prime} i}\left[\frac{\int_{j} p_{i, j, t}}{P_{i^{\prime}, t}}\right]^{-\epsilon} Y_{i^{\prime}, t}  \tag{4.23}\\
& =\left[\lambda\left(\frac{p_{i, j, t}^{*}}{P_{i, t}}\right)^{-\epsilon}+(1-\lambda)\left(1+\pi_{i, t}\right)^{\epsilon}\right] \cdot \sum_{i^{\prime}} \mu_{i^{\prime}} \gamma_{i^{\prime} i} Q_{i^{\prime}, i, t}^{\epsilon} Y_{i^{\prime}, t} \tag{4.24}
\end{align*}


Since trade linkages are a function of home bias in region 1 (denoted $\alpha$ ), we can derive the following expressions for per-capita income:

$$
\begin{aligned}
& \mathcal{Y}_{1, t}=\left[\lambda\left(\frac{p_{1, j, t}^{*}}{P_{1, t}}\right)^{-\epsilon}+(1-\lambda)\left(1+\pi_{1, t}\right)^{\epsilon}\right]\left[\alpha Y_{1, t}+(1-\alpha) Q_{2,1, t}^{\epsilon} Y_{2, t}\right] \\
& \mathcal{Y}_{2, t}=\left[\lambda\left(\frac{p_{2, j, t}^{*}}{P_{2, t}}\right)^{-\epsilon}+(1-\lambda)\left(1+\pi_{2, t}\right)^{\epsilon}\right]\left[\frac{\mu_{1}}{\mu_{2}}(1-\alpha) Q_{1,2, t}^{\epsilon} Y_{1, t}+\left(1-\frac{\mu_{1}}{\mu_{2}}(1-\alpha)\right) Y_{2, t}\right] .
\end{aligned}
$$

The above expressions are the key equations linking the trade flows between regions. Per-capita income in region $i$ is a weighted sum of regional final goods $Y_{i^{\prime}, t} \forall i^{\prime}$. If the demand for final $\operatorname{good} Y_{i^{\prime}, t}$ increases, then $\mathcal{Y}_{i, t}$ increases depending on the strength of trade linkages $\alpha$, the relative populations $\mu_{i}$, and the relative price of final good $Q_{i^{\prime}, i, t}^{\epsilon}$.

\subsection{Definition of the equilibrium}
We describe the equilibrium over the transition and leave the description of the steady-state equilibrium for Supplementary Material, Appendix F. For an exogenous sequence of regional government spending $\left\{G_{i, t}\right\}_{i=1}^{2}$, the equilibrium over the transition is a time sequence of equilibrium variables. In particular, we are looking to solve for $\left\{C_{i, t}\right\}_{i=1}^{2},\left\{L_{i, t}\right\}_{i=1}^{2},\left\{A_{i, t+1}\right\}_{i=1}^{2}$, $\left\{B_{i, t+1}^{m}\right\}_{i=1}^{2},\left\{Y_{i, t}\right\}_{i=1}^{2},\left\{\mathcal{Y}_{i, t}\right\}_{i=1}^{2},\left\{w_{i, t}\right\}_{i=1}^{2},\left\{\pi_{i, t}\right\}_{i=1}^{2},\left\{q_{i, t}\right\}_{i=1}^{2},\left\{p_{i, j, t}^{*} / P_{i, t}\right\}_{i=1}^{2}, Q_{12, t},\left\{D_{i, t}\right\}_{i=1}^{2}, R_{t}$, $\hat{\pi}_{t}, \tau_{t}$ and $\left\{\phi_{i, t}\right\}_{i=1}^{2}$, for $t=\left\{t_{0}, \infty\right\}$ where $t_{0}$ is the time of the policy change.

\begin{enumerate}
  \item Goods market equilibrium: The demand for goods by households in region $i, C_{i, t}$, is derived by the household's problem and together with local government spending, $G_{i, t}$, give the total demand for final good $i: Y_{i, t}=C_{i, t}+G_{i, t} \forall i$. The inflation rates that clear the goods market $\left\{\pi_{i, t}\right\}_{i=1}^{2}$ are derived using the following equations:
\end{enumerate}

$$
\begin{aligned}
\frac{p_{i, j, t}^{*}}{P_{i, t}} & =\frac{\epsilon}{\epsilon-1} \frac{\sum_{i^{\prime}=1}^{N} \mu_{i^{\prime}} \gamma_{i^{\prime} i} Q_{i^{\prime}, i, t}^{\epsilon}\left[w_{i, t} Y_{i^{\prime}, t}+(1-\lambda) \tilde{\beta}_{i}\left(1+\pi_{i^{\prime}, t+1}\right)^{1+\epsilon} \mathcal{X}_{i^{\prime}, i, t+1}\right]}{\sum_{i^{\prime}=1}^{N} \mu_{i^{\prime}} \gamma_{i^{\prime} i} Q_{i^{\prime}, i, t}^{\epsilon}\left[Y_{i^{\prime}, t}+(1-\lambda) \tilde{\beta}_{i}\left(1+\pi_{i^{\prime}, t+1}\right)^{\epsilon} \mathcal{Z}_{i^{\prime}, t+1}\right]} \\
1 & =\sum_{i^{\prime}} \gamma_{i i^{\prime}} Q_{i^{\prime}, i, t}^{1-\epsilon}\left[\lambda\left(\frac{p_{i^{\prime}, j, t}^{*}}{P_{i^{\prime}, t}}\right)^{1-\epsilon}+(1-\lambda)\left(1+\pi_{i^{\prime}, t}\right)^{\epsilon-1}\right] \forall i
\end{aligned}
$$

\begin{enumerate}
  \setcounter{enumi}{1}
  \item Regional income in region $i, \mu_{i} \mathcal{Y}_{i, t}$, is a weighted sum of regional final goods:
\end{enumerate}

$$
\mu_{i} \mathcal{Y}_{i, t}=\left[\lambda\left(\frac{p_{i, j, t}^{*}}{P_{i, t}}\right)^{-\epsilon}+(1-\lambda)\left(1+\pi_{i, t}\right)^{\epsilon}\right] \cdot \sum_{i^{\prime}} \mu_{i^{\prime}} \gamma_{i^{\prime} i} Q_{i^{\prime}, i, t}^{\epsilon} Y_{i^{\prime}, t} \forall i
$$

\begin{enumerate}
  \setcounter{enumi}{2}
  \item Labour market equilibrium: The labour supply satisfies the household's problem and the aggregate labour supply in region $i$ is $\mu_{i} \int_{\phi_{i, t}} x_{t} h_{t}$. Since we have a linear technology, the aggregate labour demand $\mu_{i} L_{i, t}$ equals aggregate income $\mu_{i} \mathcal{Y}_{i, t}$. The wage rate $w_{i, t}$ that clears the labour market in region $i$ is found using the following labour market condition:
\end{enumerate}

$$
\mu_{i} L_{i, t}=\mu_{i} \int_{\phi_{i, t}} x_{t} h_{t}
$$

\begin{enumerate}
  \setcounter{enumi}{3}
  \item Real exchange rate $Q_{1,2, t}=P_{1, t} / P_{2, t}$ satisfies the following equation:
\end{enumerate}

$$
\frac{\left(1+\pi_{1, t}\right)}{\left(1+\pi_{2, t}\right)}=\frac{P_{2, t-1}}{P_{1, t-1}} \frac{P_{1, t}}{P_{2, t}}=Q_{2,1, t-1} Q_{1,2, t}
$$

\begin{enumerate}
  \setcounter{enumi}{4}
  \item Dividends are given by
\end{enumerate}

$$
D_{i, t}=\left[\lambda\left(\frac{p_{i, j, t}^{*}}{P_{i, t}}\right)^{1-\epsilon}+(1-\lambda)\left(1+\pi_{i, t}\right)^{\epsilon-1}\right] \cdot \sum_{i^{\prime}} \mu_{i^{\prime}} \gamma_{i^{\prime} i} Q_{i^{\prime}, i, t}^{\epsilon} Y_{i^{\prime}, t}-w_{i, t} L_{i, t} \forall i
$$

\begin{enumerate}
  \setcounter{enumi}{5}
  \item Household savings equal the demand for government bonds and equity shares by the mutual funds
\end{enumerate}

$$
\sum_{i} \mu_{i} A_{i, t+1}=\sum_{i} \mu_{i}\left(1+\pi_{i, t+1}\right) B_{i, t+1}^{m}+\sum_{i} \mu_{i} q_{i, t},
$$

where the total number of shares is normalized to one and the value of equity, $q_{i, t}$, is given by arbitrage equation (4.9).\\
7. The demand for government bonds by the mutual funds equals the supply of bonds by the government: $\sum_{i} \mu_{i} B_{i, t}^{m}=B_{t}^{g}$. The supply of bonds along the transition is given by the fiscal rule: $B_{t}^{g}=B_{s s}^{g}-\sum_{i} \mu_{i}\left(q_{i, t}-q_{i, s s}\right)$.\\
8. The tax rate $\tau$ is found by balancing the government budget constraint:

$$
\sum_{i} \mu_{i}\left(1+\pi_{i, t+1}\right) B_{i, t+1}^{m}-\left(1+R_{t-1}\right) \sum_{i} \mu_{i} B_{i, t}^{m}=\sum_{i} \mu_{i} G_{i, t}-\sum_{i} \mu_{i} \int_{\phi_{i, t}} \mathcal{T}_{t}-\tau_{d} \omega \sum_{i} \mu_{i} D_{i t} .
$$

\begin{enumerate}
  \setcounter{enumi}{8}
  \item The monetary policy does not respond to inflation: $R_{t}=R_{s s}$.
  \item The national inflation rate is given by: $\hat{\pi}_{t}=\sum_{i=1}^{2} \mu_{i} \pi_{i, t}$.
  \item The regional measures $\phi_{i, t}$ evolve based on the policy functions and the transition matrices described in the model.
\end{enumerate}

For the steady-state equilibrium, we assume inflation is zero, and prices are symmetric within and across regions: $p_{i, j} / P_{i}=1 \forall i$ and $Q_{i i^{\prime}}=1 \forall i, i^{\prime}$.

\subsection{Existence and uniqueness of equilibrium}
We impose two restrictions when computing our equilibrium.\footnote{The computational algorithm to solve for the transitional dynamics is presented in Supplementary Material, Appendix I.} First, in our equilibrium, we restrict attention to transition paths for which at a sufficiently long time after the fiscal stimulus is over, the inflation rate in both regions is equal to its steady-state value, i.e. $\pi_{i, t^{*}}=0$, for all regions, and the relative price is equal to its steady-state value, i.e. $Q_{12, t^{*}}=1$, where $t^{*}$ is sufficiently large. As such, equilibrium paths converge over time to the initial steady state. Since the monetary policy (i.e. the nominal interest rate) does not respond to inflation, it is possible that there are multiple equilibrium paths that lead to the original steady state.

As a result, we additionally impose restrictions to the state space. Specifically, we focus on the minimum state variable (MSV) equilibrium which arises when the optimal choices depend on a minimal set of state variables.\footnote{One motivation for restricting attention to MSV equilibria appears in Angeletos and Lian (2021). They show that an infinitesimally small memory loss on behalf of future agents renders all but one equilibrium, the MSV equilibrium, explosive. Note, however, that Angeletos and Lian (2021) study linear models.} Our set of state variables is selected such that it is not possible to delete a single state variable (or a group of state variables) from agents' policy functions and continue to obtain a solution that satisfies the model's equilibrium conditions (McCallum, 1983). We could have selected extraneous state variables (e.g. sunspot shocks) whose inclusion could potentially give another solution satisfying the equilibrium conditions but would not be necessary to obtain one. Focusing on the class of MSV equilibria, the MSV solution is the unique equilibrium in linear settings (McCallum, 1983). We do not know with certainty if uniqueness generalizes to nonlinear dynamic systems as in our model.\footnote{Hagedorn~\cite{Hagedorn2016-un} and Hagedorn et al.~\cite{Hagedorn2019-op} offer an alternative possibility to avoid multiple equilibria. In their framework, uniqueness of equilibrium is guaranteed by a combination of incomplete markets (and thus, precautionary savings) and nominal bond targeting. This result does not apply to our case where we make the more standard assumption that the government targets real debt.}


% Smart bibliography: Only include bibliography if standalone AND has citations
\smartbib

\pdfonly{
  % execute this version if compiling with pdflatex
  \captionsetup[figure]{list=no}
  \captionsetup[table]{list=no}
}

\end{document}
